% 本文件由 LaTeX 排版 Agent 根据有效 translation.json 自动生成。
% author=Dionysius the Great; work_id=0613; title=释经残篇（狄约尼削）

\chapter{释经残篇（\Person{狄约尼削}）}

% structure: p0001 source=h1 target=omitted reason=page-title-already-rendered-by-parent

% structure: p0002 source=h2 target=section reason=relative-heading-rank; base=h2

\section{论\BibleRef{路 22:42-48}}

\BibleQuote{父啊，祢若愿意，请把这杯从我面前撤去；但不要照我的意愿，而要照祢的意愿成就。}\par

但关于意志，说这些就够了。然而\Quote{让这杯过去}这句话，并不意味着\Quote{让它不要靠近我}或\Quote{接近我}。因为能从祂那里过去的事物，必定先靠近祂；而这样从祂那里过去的事物，也必须经过祂。因为如果它没有到达祂，就不能从祂过去。因为祂承担了人的位格，作为成了人的人。因此，在此场合，祂反对办理那较低者，即祂自己的意愿，而祈求那较高者，即祂圣父的意愿得以成就，也就是天主的旨意；而这旨意，在神性方面，在祂自己与圣父内原是同一的。因为圣父的旨意是要祂经历各种试探；圣父以奇妙的方式带领祂走上这条路，并非以试探本身为目标，也不是为了简单进入试探，而是为了证明祂高于试探，也超越试探。确实，救主没有祈求不可能的事，也没有祈求不可行的事，也没有祈求违反圣父旨意的事。这是可能的；因为\Person{马尔谷}提到祂说：\BibleQuote{阿爸，父啊！一切为祢都可能。}而且如果祂愿意，它们就是可能的；因为\Person{路加}告诉我们祂说：\BibleQuote{父啊，祢若愿意，请把这杯从我面前撤去。}因此，圣神在圣史们之间分配，通过这几位叙述者共同的话语，构成了救主整个心意的完整记述。所以，祂并没有向圣父祈求圣父所不愿意的事。因为\Quote{你若愿意}这句话，表明的是服从和顺从，而非无知或犹豫。为此，另一段圣经说：\BibleQuote{一切为祢都可能。}\Person{玛窦}也极好地描述了服从和谦卑，当祂说：\BibleQuote{若是可能。}因为除非我这样调整理解，否则有人可能会给\Quote{若是可能}这个表达赋予不虔敬的意义，好像天主有什么不可能做的事，除了祂不愿意做的事。但是……祂的人性被祂固有的神性立即坚强起来，于是祂提出更安全的祈求，不再愿意那样，而是愿它按照圣父的善意、在荣耀、恒心和圆满中完成。因为\Person{若望}——他给我们记录了救主最高深、最神圣的言行——听见祂这样说：\BibleQuote{圣父赐给我的杯，我岂能不喝吗？}现在，喝这杯就是勇敢地承担那职务和整个试探的经济，追随并完成圣父的决定，并克服一切恐惧。而\Quote{祢为什么舍弃了我？}的呼喊，与祂先前所作的祈求完全一致：\BibleQuote{为什么死亡一直与我同在直到现在，而我还没有承受这杯？}我认为这就是救主在这简练表达中的意思。\par

祂那时说真话。但祂并没有被舍弃。然而祂立刻饮尽了这杯，正如祂所祈求的，然后逝去。给祂递来的醋，在我看来是一件象征性的事。因为变酸的酒很好地表明了祂所承受的迅速转变与改变，当祂由受难转变为不受苦难，由死亡转变为不死，由受审者的地位转变为审判者，由受暴君权下转变为行使王权。而海绵，我想，表示了祂内实现的圣神的完全灌注。芦苇象征了王权与神圣法律。牛膝草表达祂那赋予生命与救赎的复活，藉此祂也给我们带来了健康。\par

\BibleQuote{有一位天使从天上显现给祂，加强祂的力量。}\par

\BibleQuote{祂在极度痛苦中，祈祷更加恳切；祂的汗如同大血点滴在地上。}\par

\Quote{血汗}这句话是当前比喻性说法，用于极度痛苦和忧愁的人；正如有人极其悲伤，人们说那人流泪如血。因为用了\Quote{如同大血滴}的表达，他并没有说汗滴实际上就是血滴。因为若是这样，他就不会说这些汗滴像血了。因为\Quote{如同大滴}就是这样的力量。反而，为了表明主的身体不是被任何细微的、只有表象的湿气所沾湿，而是全身都被大而浓的汗珠所湿透，他以大血滴来比喻祂的情况。因此，通过恳求的强度与剧烈的痛苦，也通过浓密而大量的汗，他揭示了这一事实：救主在本质上和实际上是人，而非仅在外表与形象上，并且祂承担了人性中所有无辜的情感。然而，\BibleQuote{我有权舍掉我的生命，也有权再取回它}这句话，表明祂的受难是自愿的；此外，它指出被舍去和取回的生命是一回事，而舍去和取回它的神性是另一回事。\par

祂说\Quote{一回事与另一回事}，不是将祂分作两个位格，而是显示两性之间的区别。\par

再者，正如祂甘愿在肉身承受死亡，从而将不朽植入肉身；同样，祂凭自己的自由意志承担我们奴役的痛苦，从而在其中播下恒毅和勇气的种子，由此使那些信祂的人得以坚强，为他们作见证的剧烈战斗作好准备。因此，那些汗滴也以奇妙的方式从祂身上流下，如同大血滴一般，为的是要，可以说，排干并倒空本性中固有的恐惧之泉。因为若非此举具有奥义，即使祂是最胆怯、最卑微的人，也绝不会在祂的痛苦压力下，以这种非自然的方式被汗滴如同血滴般浸湿。\par

叙述中告知我们有一位天使站在救主身旁并坚固祂的那句话，也具有同样意义。因为这也关系到为我们所实行的救恩计划。那些被委派为虔诚而进行神圣战斗的人，自有天上的天使前来协助他们。至于那句祈祷：\Quote{父啊，请撤去这杯}，祂这样说大概并非因为祂怕死本身，而是要用这些话挑战魔鬼，让魔鬼为祂竖立十字架。那一位曾以欺骗之言迷惑了亚当；现在，愿那欺骗者自身以神性之言被迷惑。然而，圣子的旨意与圣父的旨意并非不同。因为凡是愿意圣父所愿意的，就被发现具有圣父的旨意。所以祂说\Quote{不要照我的意愿，唯照祢的意愿}，这是一种比喻的说法。并非祂希望那杯被撤去，而是将祂苦难的正确结局交托于圣父的旨意，并以此尊崇圣父为元首。因为如果教父们将人的性情称为 \OriginalTerm{意志}{gnomè}，而这样的性情也涉及如同定意之中隐藏的考量，那么有些人怎能说那超乎这一切的主具有一种 \OriginalTerm{意志的意愿}{gnomic will}？显然，这只能是由于理性的缺失。\par

45. 祂从祈祷中起来，来到门徒那里，见他们因忧愁都睡着了；\par

46. 就对他们说：\Quote{你们为什么睡觉呢？起来祈祷吧，免得陷于诱惑。}\par

因为最普遍地说来，似乎没有人能完全免于邪恶的经验。正如有人所说：整个世界都卧于邪恶之中；又说：人的日子大多是劳苦和烦恼。但你或许会说：受试探与跌倒或陷于诱惑之间有什么不同？嗯，如果一个人被邪恶战胜——而且除非他自己与之斗争，除非天主以祂的盾牌保护他，他必被战胜——那么这个人就陷于诱惑，在里面，并如同被俘虏一般被压制。但如果一个人抵抗并忍受，那么他确实受了试探，却没有陷于诱惑或跌入其中。因此，耶稣曾被圣神带领，不是去陷于诱惑，而是去受魔鬼试探。同样，亚巴郎并没有陷于诱惑，天主也没有领他进入诱惑，而是试探（考验）了他；但并没有驱使他进入诱惑。主自己也试探（考验）门徒。因此，那恶者试探我们时，是将我们拖入诱惑，因为祂本身是以邪恶的诱惑行事。但天主试探（考验）时，是作为不受邪恶试探者提出考验。因为经上说：天主不能被邪恶试探。所以，魔鬼强暴地驱使我们，将我们引向灭亡；而天主引导我们，为我们的救恩而训练我们。\par

47. 祂正说话的时候，看啊，来了一群人，那称为犹达斯的，十二人中之一，走在前头，接近耶稣，并亲吻了祂。\par

48. 耶稣对他说：\Quote{犹达斯，你用亲吻来出卖人子吗？}\par

主的容忍是多么奇妙！祂甚至亲吻了那出卖者，并且说出比亲吻更温和的话！因为祂并没有说：\Quote{你这可憎的，甚至极其可憎的叛徒，你就是这样回报我们如此巨大的恩惠吗？}而是简单地说了\Quote{犹达斯}，使用了专名，这是一种怜悯人或想要将人唤回者的称呼，而非愤怒者的语气。祂也没有说\Quote{你的师傅、主、恩人}，而只是说\Quote{人子}，即温柔和温良的那一位：仿佛祂的意思是：\Quote{即使我不是你的师傅、主或恩人，你仍要出卖一位对你如此无伪、如此温柔的人，甚至在你叛逆的时刻亲吻你，而且这亲吻还是你背叛的信号吗？}主啊，祢当受赞美！祢亲自向我们显示了何等伟大的容忍榜样！何等谦卑的楷模！然而，主给了我们这个榜样，为要向我们表明，即使我们的话语没有产生显著效果，我们也不应放弃向弟兄们提供善意的劝告。\par

因为正如不愈之伤是指那些无论以严厉的方法，还是以较为柔和的方法都无法治愈的伤口；同样，灵魂一旦被掳，并自甘堕落于任何一种邪恶，且拒绝考虑真正有益于它的事物时，即使有万千劝告在其内回响，它对自己也毫无益处。正如听觉在其中死亡，它从向它发出的劝诫中得不到益处；不是因为不能，只是因为不愿。这就是发生在犹达斯身上的事。然而基督，虽然事先知道这一切，但从始至终，从未停止做一切凭祂所能的劝告之事。既然我们知道祂是如此行事，我们也当不懈地努力纠正那些漫不经心的人，即使这些劝告似乎没有产生实际的益处。\par

% structure: p0019 source=h2 target=section reason=relative-heading-rank; base=h2

\section{论\BibleRef{路 22:42} 等}

但关于这旨意，说这些就够了。然而，\BibleQuote{让这杯离去}这句话的意思，并不是说\Quote{不要它临近我}或\Quote{不要它靠近我}。因为凡能从祂离去的，必定先临近祂；而确实从祂离去的，也必定由祂而出。因为若没有临到祂，就不能从祂离去。因此，如同祂现在感到它就在眼前，祂就开始忧闷、惊慌、极为恐惧，并陷入苦痛。又如同它已临近并摆放在祂面前，祂不只说\Quote{杯}，而是用\Quote{这}字来指明。所以，正如从一人离去之物，既不是毫无临近，也不是永远停留，同样，救主的第一个祈求是，那轻轻悄悄临到祂、轻轻与祂结合在一起的试探可以被移开。这就是祂也教导软弱的门徒们祈求的那种不陷入试探的第一种形式；那涉及试探的临近：因为绊倒的事是免不了的，但那些遭遇它们的人不应陷入试探。但祂在第二个祈求中所表达的不进入试探的最圆满方式，是当祂不仅说\BibleQuote{不要照我的意愿}，而且说\BibleQuote{但要照祢的意愿}时。因为在天主那里没有恶的试探；祂愿意赐给我们超过我们所求所想的丰盛恩惠。因此，祂的旨意是圆满的旨意，爱子自己知道；祂常常说祂来是为承行那旨意，而非祂自己的旨意——即人的旨意。因为祂取了人的位格，成为人。因此，此时祂也祈求那较低的、属于祂自己的旨意不被实行，而祈求那较高的、属于祂圣父的旨意得以实行，即天主的旨意，而这在神性方面，在祂自己和圣父内是同一旨意。因为圣父愿意祂经历一切试探，并且圣父以奇妙的方式引领祂走上这条路；的确，不是以试探本身为目的，也不是仅仅为了让祂进入试探，而是为了证明祂超越试探，也超脱试探。确实，救主所祈求的既不是不可能的事，也不是不可行的事，也不是违反圣父旨意的事。这是可能的事，因为马尔谷提到祂说：\BibleQuote{阿爸，父啊，一切为祢都是可能的}；如果祂愿意，这些都是可能的，因为路加告诉我们祂说：\BibleQuote{父啊，祢若愿意，请把这杯从我这里撤去}。因此，圣神分施在福音作者们中间，通过他们几位叙述者的话语共同构成了我们救主全部心境的完整记述。所以，祂并没有向圣父祈求圣父所不愿的事。因为\BibleQuote{祢若愿意}这话显示的是服从和温顺，而非无知或犹豫。正如当我们向父亲、统治者或我们所尊敬的人提出合乎他们判断的请求时，我们惯常使用\Quote{若祢喜欢}这样的称呼（当然并不是因为我们对此有怀疑）；同样，救主也说\BibleQuote{祢若愿意}——并非祂以为祂愿意别的事，然后才得知实情，而是祂确切知道天主愿意将杯从祂撤去，并且如此行时也正确理解祂所愿意的也是为祂可能的。因此另一部经上说：\BibleQuote{一切为祢都是可能的}。玛窦又极好地描述了这种服从和谦卑，当他说\BibleQuote{若是可能}。因为除非我们这样理解，否则有人或许会赋予\BibleQuote{若是可能}这句话一个不敬的意义，仿佛天主有什么事是不可能的，除非是祂不愿意做的事。因此祂所作的祈求并非独立自足的，也不是只取悦祂自己或反对圣父旨意的，而是符合天主心意的。然而，有人可能会说祂被压倒而改变心意，随后祈求与先前不同的事，不再坚持自己的意愿，而是引入祂圣父的意愿。的确，情况正是如此。但祂并非从一方转向另一方；而是采用另一种方式、另一种方法来完成同一件事，这也是双方都喜欢的；即，选择取代那较低的、似乎也令祂自己不满的方式，而采纳圣父所指示的更高、更令人钦佩的方式。因为无疑祂曾祈求杯能离去；但祂也说：\BibleQuote{然而，不要照我的意愿，但要照祢的意愿}。祂一方面痛苦地切望它离去，但（祂知道）照圣父的意愿更好。因为祂并非祈求它现在不离去以替代撤去；而是当撤去已摆在眼前时，祂宁愿这按圣父的意愿成就。因为撤去有两种：一种是在对象显现并触及另一者之后，随即被追赶或超越而离去，如同赛跑者彼此擦身而过；另一种是对象在另一者处逗留、停留并坐下，如同掠夺队或营帐，过段时间因被征服而撤去，且并未得胜。因为若他们得胜就不撤退，而是带走他们所征服的人；若他们不能得胜，就耻辱地自行撤退。祂是按照第一种样式希望杯来到祂手中，并迅速轻易地离去；但祂一说这话，立刻因圣父的神性在人性中得到加强，祂就提出更稳妥的祈求，不再希望那样，而是希望按圣父的美意、在光荣、坚定和圆满中成就。因为若望（他记录了救主最崇高和最神圣的言行）听见祂这样说：\BibleQuote{父所给我的杯，我岂可不喝呢？}而喝这杯，就是以坚毅完成使命和整个试探的经纶，追随并实现圣父的决定，并超越一切恐惧；确实，在祂所作的祈祷中，祂表明自己已将这些恐惧抛在身后。因为两个对象中，任何一个都可以说从另一个移开：留存的对象可以说从离开的那一个移开，而离开的那一个也可以说从留存的那一个移开。此外，玛窦最清楚地表明，祂确实祈求杯能离去，但祂的祈求是这应不照祂的意愿，而照圣父的意愿成就。马尔谷和路加所记的话也应在适当之处引出。因为马尔谷说：\BibleQuote{然而不要照我所愿意的，但要照祢所愿意的}；路加说：\BibleQuote{然而不要照我的意愿，但要照祢的意愿成就}。祂确实表达了那样的意思，并确实渴望祂的苦难迅速减退并结束。但与此同时，圣父愿意祂以持久努力和足够充分的方式进行战斗。因此，圣父（祂）引入了所有攻击祂的事物。但射向祂的标枪中，有些被击碎，有些被如同无敌的钢铁臂膀弹回，或者更像是从坚固不动的磐石上弹回。打击、唾弃、鞭笞、死亡以及死亡中的高举，这一切都临到祂；当这一切都经过后，祂沉默忍耐到底，如同什么都没受，或已死了一样。但当祂的死亡被拖延，且（如果可以这样说）超过了祂的极限力量时，祂向圣父呼喊：\BibleQuote{祢为什么舍弃了我？}这个呼喊与祂先前所作的祈求完全一致：\Quote{为什么死亡一直如此紧密地与我同在直到现在，而祢还不将杯从我这里过去？难道我还没有喝它、喝尽它吗？但若没有，我害怕它持续的压力会完全耗尽我；而这就是若祢舍弃我时所发生的事：那时，应验会存留，但我却会消逝，变得毫无效力。现在，我恳求祢，让我的洗礼完成吧，因为我确实极其迫切直到它成就。}——我认为这就是救主在这简短话语中的意思。祂那时当然说的是真话。然而祂并未被舍弃。尽管如此，祂立刻喝尽了杯，正如祂的恳求所暗示的，然后逝去。递给祂的醋，在我看来是一件象征性的事。因为变酸的酒很好表明了祂从受难到不受难、从死亡到不死、从被审判者的地位到审判者的地位、从暴君权下的服从到行使君王统治时所经历的快速转变和更迭。而海绵，我想，代表在祂身上实现的圣神的完全灌注。芦苇象征王权和神圣的法律。牛膝草表达了祂那给予生命和拯救的复活，藉此祂也给我们带来了健康。但我们已在玛窦和若望福音上足够详细地讨论了这些事。若蒙天主许可，我们也将讨论马尔谷的记述。但现在，我们将继续讨论我们段落中接下来的内容。\par

% structure: p0021 source=h2 target=section reason=relative-heading-rank; base=h2

\section{论\BibleBook{路加}{福音}22:46等处}

这篇祈祷祂自己也献上了，屡次俯伏在地；在两次场合中，祂都恳求不要进入试探：一次祂祈祷说：\BibleQuote{如果可以，请叫这杯离开我}；另一次祂说：\BibleQuote{然而不要照我的意愿，但照祢的意愿。}因为祂说到不要进入试探，并以此为祈祷；但祂并未祈求在这些情况下全然不受磨难，或任何苦难永不临于祂。因为最普遍的应用来说，任何人完全不受恶的经验似乎不可能；正如有人所说：\BibleQuote{整个世界都卧于邪恶中}；又有人说：\Quote{人的日子多半是劳苦和烦恼}，人们自己也承认。\Quote{我们的生命短暂，且充满忧愁。}然而，祂不宜直接命令他们祈祷那诅咒不要应验，这诅咒如此表达：\BibleQuote{地因你的作为而受诅咒：你一生日日劳苦才能从地里得食}；或\BibleQuote{你本是尘土，仍要归于尘土。}为此，圣经在许多方面以不同方式指明生活的悲惨困苦，称之为\Quote{流泪谷}。尤其对于圣人们，这个世界更是痛苦之地，祂对他们说这话，祂说话绝不撒谎：\BibleQuote{在世界上你们将有苦难。}同样，祂也藉先知说：\BibleQuote{义人的苦难众多。}但我认为祂所指的不要进入试探，是当祂用先知的话说从苦难中被拯救时。因为祂接着说：\BibleQuote{主要救他脱离这一切。}这正符合救主的话，祂应许他们将胜过苦难，并分享祂为他们赢得的胜利。因为祂先说\BibleQuote{在世界上你们将有苦难}，然后加上\BibleQuote{但放心吧，我已经胜了世界。}此外，祂教导他们祈祷不要陷入试探，说：\BibleQuote{不要让我们陷入试探}；这意思是\BibleQuote{不要容许我们陷入试探。}为了表明这并不表示他们不应受试探，而是实在要脱离邪恶，祂加上\BibleQuote{但救我们脱离凶恶。}但你或许会说：\Quote{受试探与陷入或进入试探有何区别？}嗯，若一个人被恶所胜——除非他自己抵抗，且天主用盾牌保护他，否则他必被胜——那人便已进入试探，且在试探中，被挟制如同被掳者。但若有人抵挡并忍耐，那人确实受了试探，但并未进入或陷入试探。因此耶稣被圣神领往旷野，并非要进入试探，而是受魔鬼的试探。同样，亚巴郎并未进入试探，天主也没有领他进入试探，而是试探（试炼）了他；但并未驱使他进入试探。主自己也试探（试炼）了门徒。因此恶者试探我们时，使我们陷入试探，因为他亲自处理恶的试探；但天主试探（试炼）时，祂是作为不受恶试探者带来试探。因为经上说：\BibleQuote{天主不能被恶试探。}所以魔鬼强行驱使我们走向毁灭；但天主引导我们的手，训练我们走向救恩。\par

% structure: p0023 source=h2 target=section reason=relative-heading-rank; base=h2

\section{\BibleBook{}{训道篇}开篇注释}

% structure: p0024 source=h3 target=subsection reason=relative-heading-rank; base=h2

\subsection{第一章}

1节。\BibleQuote{\Emph{言语}是达味之子、在耶路撒冷为王的言语。}\par

同样，玛窦也称主为达味之子。\BibleRef{玛 1:1}\par

3节。\BibleQuote{人在太阳之下劳碌，一切劳碌中有什么益处？}\par

因为有什么人，虽因劳碌追求地上的事物而变得富有，却能将自己由二肘身高变为三肘呢？或有什么人，原为瞎子，却以此恢复了视力？因此我们应将劳碌指向超越太阳的目标：因为德行同样达到那里。\par

4节。\BibleQuote{一代过去，一代又来，大地却永远长存。}\par

是的，长存（到世代），但不是到永世。\par

16节。我心里议论说：看哪，我已大有智慧，比我以前在耶路撒冷的众人更有智慧；我的心也多经历智慧和知识。\par

17节。我明白箴言和科学：知道这也是精神的抉择。\par

18节。因为多有智慧，就多有知识；加增知识的，就加增忧伤。\par

我徒然自高自大，增添了智慧；并非天主所赐的智慧，而是\BibleQuote{这世界的智慧在天主面前是愚拙的}的这种智慧\BibleRef{格前 3:19}。因为在这方面，撒罗满的经验也超越了审慎，并超过所有古人的尺度。因此他显明了其虚妄，正如下文所示：\BibleQuote{我的心说出许多事：我认识智慧、知识、箴言和科学。}但这并非真正的智慧或知识，而是如保禄所说，使人自高自大的那种。他又如经上记载\BibleRef{列上 4:32}说了三千则箴言。但这些并非属灵的箴言，而只是适合一般人类社会生活的；例如关于动物或医药的言论。因此他以戏谑的口吻加上了\BibleQuote{我知道这也是精神的抉择。}他也说到知识的增多，并非圣神的知识，而是此世王子所运作的，并用以诱骗人的灵魂，以好奇的问题追问天的高度、地的位置、海的界限。但他也说\BibleQuote{加增知识的，就加增忧伤。}因为他们甚至探求更深的事——例如，询问火为何必须上升，水为何必须下降；当他们得知是因为一轻一重时，只是增加了忧伤：因为问题仍在：为何不能相反呢？\par

% structure: p0035 source=h3 target=subsection reason=relative-heading-rank; base=h2

\subsection{第二章}

第1节。\BibleQuote{我在心里说：来吧，以欢乐试试，并在幸福中观看。看，这也是空虚。}\par

因为他是为了试探，并依照高超而严格的生活所带来的，才进入享乐。他提到人们所谓的欢乐。他说\Quote{在幸福中}，指的是人们所谓的福乐，这些福乐不能给拥有者生命，并使沉浸在其中的人如同它们一样空虚。\par

2. \BibleQuote{我论欢笑说：这是疯狂；论欢乐说：你在做什么？}\par

欢笑具有双重的疯狂：因为疯狂产生欢笑，却不允许为罪忧伤；又因为那样的人被疯狂所附，混淆时节、地点和人物。因为他逃避那些忧伤的人。\BibleQuote{论欢乐说：你在做什么？}你为何到那些不能欢乐的人那里去？为何到醉酒者、贪婪者和掠夺者那里去？何以这措辞如同酒？因为酒使人心欢乐；它对贫穷困苦的人有作用。然而，肉体当以规律和适度行动时，也使人心欢乐。\par

3. \BibleQuote{我的心在智慧中引导我，并在欢乐中得胜，直到我知道什么是人子们当在太阳下，在他们一生的天数中所行的善事。}\par

他说，受智慧引导，我在欢乐中战胜了享乐。而且，我的知识目标是专注于不空虚的事，而是寻找善；因为人若找到善，也不会错过对有益之事的辨识。充足也是适时，与生命的长度相称。\par

4. \BibleQuote{我为自己作了大工程，建造了房屋，栽种了葡萄园。}\par

5. \BibleQuote{我为自己造了花园和果园。}\par

6. \BibleQuote{我为自己造了水池，用以浇灌生长树木的林园。}\par

7. \BibleQuote{我买了仆婢，也有生在家中的仆婢；此外，我拥有大群和小群的牲畜，多过耶路撒冷以前所有的人。}\par

8. \BibleQuote{我又为自己积蓄金银，以及君王和各省的珍宝。我为自己召集了歌唱的男女，以及人子们的欢乐，如同杯和持杯者。}\par

9. \BibleQuote{于是我成了伟大，超过耶路撒冷以前所有的人；并且我的智慧仍与我同在。}\par

10. \BibleQuote{凡我眼所愿的，我未加禁止；我不拒绝我的心享有一切欢乐。}\par

你看见他如何列举众多房屋、田地以及他所提及的其他事物，然后发现其中毫无益处。因为他不因这些事物在灵魂上变得更好，也未借此获得与天主的友谊。他必然也被引导谈及真财富和恒久的产业。因此，他想显示哪种财产会留在拥有者手中，稳固不变并为他存留，他加上：\BibleQuote{并且我的智慧仍与我同在。}因为唯独这智慧存留，而他已经列举的其他一切都逃逸消失。因此，智慧仍与我同在，而我因它而存留。因为那些事物跌倒，也使追逐它们的人跌倒。但是，为了在智慧与人们视为善的事物之间作比较，他加上这些话：\BibleQuote{凡我眼所愿的，我未加禁止}等等；借此他描述为恶的，不仅是那些为满足自己享乐而劳苦的人所忍受的辛劳，还有那些为维系日常生活而必须和被迫承受的辛劳，每日汗流满面从事不同的劳作。他说，劳苦是大的，但借着劳苦的技艺是暂时的，在令人喜悦的事物中增添无益之物。因此没有益处。因为哪里没有卓越，哪里就没有益处。所以，这些忧虑的对象有理由成为空虚和\OriginalTerm{神魂}{spirit}的选择。\OriginalTerm{神魂}{spirit}这个名，他用来指\OriginalTerm{灵魂}{soul}。因为选择是一种品质，而非运动。达味说：\BibleQuote{我将我的\OriginalTerm{神魂}{spirit}交在祢手中。}确实，\BibleQuote{我的智慧仍与我同在}，因为它使我认识并理解，以便我能讲述太阳下一切无益之事。因此，如果我们渴望公义的益处，如果我们寻求真正的利益，如果我们以不朽为目标，让我们从事那些超越太阳的劳作。因为在其中没有空虚，也没有对一个既空虚又无目的地匆忙来去的\OriginalTerm{神魂}{spirit}的选择。\par

12. \BibleQuote{我转而观看智慧、疯狂和愚昧：因为那将随从智慧在其所行的一切事上之后而来的人是谁呢？}\par

他指来自天主的智慧，这智慧也与他同在。他借着疯狂和愚昧指称人的一切劳作，以及其中空虚而愚蠢的快乐。因此，区分这些及其尺度，并祝福真智慧，他加上：\BibleQuote{因为那将随从智慧之后而来的人是谁呢？}因为这智慧教导我们那真正的智慧，并赐予我们脱离疯狂和愚昧的拯救。\par

13. \BibleQuote{于是我看出智慧胜过愚昧，如同光明胜过黑暗。}\par

他这样说并非出于比较。因为彼此相反、互相毁灭的事物不能比较。但他的决断是，一个应当选择，另一个应当避免。类似的说法是：\BibleQuote{世人\OriginalTerm{宁可}{rather}爱黑暗而不爱光明。}\BibleRef{若 3:19}因为那段经文中的\Quote{\OriginalTerm{宁可}{rather}}一词表达了爱好者的选择，而非对象本身的比较。\par

14. \BibleQuote{智慧人的眼目在他头上，愚昧人却在黑暗中行走。}\par

他说，人总是倾向于属地之事，并且统治官能变得昏暗。确实，如果我们只谈论身体的构造，我们人类都有眼睛在头上。但他这里说的是心灵的眼睛。因为就如猪的眼睛不会自然地朝向上天，正是因为它的天性倾向于肚腹；同样，一个曾被享乐削弱之人的心灵，也不容易从已经养成的倾向中转向，因为他没有\Quote{遵守主的一切诫命}。再者：基督是教会的头。\BibleQuote{基督是教会的头。}\BibleRef{弗 5:23}因此，那些行走在祂道路上的人就是智慧的人；因为祂亲自说过：\BibleQuote{我就是道路。}\BibleRef{若 14:6}因此，智慧的人应当总是将心灵的眼睛注视基督本身，以便他行事适度，不在顺境中心高气傲，也不在逆境中疏忽懈怠：\Quote{因为祂的判断如同深渊，}正如你将从下文更确切地了解。\par

14. 我也看出自己，同一件事临到他们众人。\par

15. 于是我在心里说：愚昧人所遭遇的，我也遭遇了，那么我为什么更有智慧呢？\par

接下来论述的进程针对那些对此生持卑下看法的人，在他们看来，死亡之事和身体的一切反常痛苦都是一种可怕的灾祸，并因此认为修德生活毫无益处，因为在这些祸患中，智慧的人与愚蠢的人没有区别。因此，他将这些视为近乎完全疯狂的言语；因此他加上这句话：\Quote{因为愚蠢的人说话过多；}这里的\Quote{愚蠢的人}指的是他自己，以及所有那样推理的人。因此他谴责这种荒谬的思维方式。出于同样的原因，他在内心的恐惧中发出了这样的感慨；并且害怕要听的人正直的谴责，他通过自己的反思解决了迫切的难题。因为这句话：\Quote{为什么我当时更智慧？}是一个处于怀疑和困难中的人的话，他怀疑用于智慧上的努力是否做得对或徒劳无功，以及智慧的人与愚蠢的人在利益上是否有区别，既然前者与后者同样受到今世发生的相同苦难。因此他说：\Quote{我在心里说得太过分了，}认为智慧的人与愚蠢的人没有区别。\par

16. \Quote{因为智慧的人和愚昧的人一样，永远无人记念。}\par

因为今世发生的一切事件都是短暂的，即使是痛苦的遭遇，如他所说：\Quote{现在一切都被交予遗忘。}因为经过短暂的时间，今世发生在人身上的一切都消失在遗忘中。是的，经历过这些事的人并不都被同样记念，即使他们经历过生活中相似的机遇。因为他们被记念不是因为这些，而是因为他们所表现出的智慧或愚昧、德行或恶行。这些人的记忆不会因临到他们命运的变迁而在人中间（同样）消灭。因此他加上：\BibleQuote{智慧的人怎能与愚昧的人一同死去？罪人的死亡确实是恶的；然而义人的记忆蒙福，恶人的名字却必朽坏。}\BibleRef{箴 10:7}\par

22. \Quote{因为那临到人在他一切劳碌上。}\par

实话实说，对于那些心思被世俗分心占据的人，生命成为一件痛苦的事，它仿佛用它的刺刺伤人心，也就是用贪求更多的情欲。贪婪所带来的忧虑也是可悲的：它并不怎么使成功的人满足，反而使不成功的人痛苦；白天在劳苦的焦虑中度过，黑夜因求利的挂虑使睡眠从眼前逃走。因此，那注视这些事物之人的热忱是虚妄的。\par

24. 人除了吃、喝，并在劳碌中使灵魂得享美善，别无好处。这也我看见，是出于天主的手。\par

25. 因为谁从自己的资源中吃喝呢？这段论述现在不涉及物质的食物，他将通过以下内容表明：即\BibleQuote{往哀悼的家去，胜过往宴乐的家去。}\BibleRef{训 7:2}因此，在当前的段落中，他继续添加：\Quote{并在劳碌中使（什么）显为灵魂的美善。}当然，仅仅是物质的食物和饮料并不是灵魂的美善。因为肉体在奢华滋养下，会与灵魂交战，并起来反抗精神。狂饮暴食岂不也违背天主吗？因此，他谈论的是奥秘之事。因为没有人能分享这属灵的筵席，除非他被祂所召，并听从那说\BibleQuote{拿起来吃}的智慧。\BibleRef{箴 9:5}\par

% structure: p0065 source=h3 target=subsection reason=relative-heading-rank; base=h2

\subsection{第三章}

3. \Quote{有时杀戮，有时医治。}\par

对于\Quote{杀戮}，是指那些犯下不可赦免之过犯的人；对于\Quote{医治}，是指那些呈现可医治伤口的人。\par

4. \Quote{有时哭泣，有时欢笑。}\par

哭泣之时，就是受苦之时；正如主也说：\BibleQuote{我实在告诉你们，你们将要哭泣哀号。}\BibleRef{路 6:25}而欢笑，则关乎复活：\Quote{因为你们的忧愁，}祂说，\BibleQuote{要变为喜乐。}\BibleRef{若 16:20}\par

4. \Quote{有时哀悼，有时舞蹈。}\par

当人想起亚当的过犯带给我们的死亡时，就是哀悼之时；但当我们记念我们期待通过新亚当而来的死人复活时，就是举行庆节之时。\par

6. \Quote{有时保留，有时丢弃。}\par

有保存圣经而不令不配者见之时，亦有将其交付堪当者之时。或者，再者：在降生成人以前，是保存法律字面之时；但当真理开花时，便是抛弃它之时。\par

7. \BibleQuote{有缄默之时，亦有发言之时。}\par

有发言之时，即当听者接纳言语之时；亦有缄默之时，即当听者曲解言语之时；正如保禄所说：\BibleQuote{异端者，在第一次和第二次劝戒以后，应拒绝他。} \BibleRef{铎 3:10}\par

10. 于是，我看见了天主赐予人类子孙在其中操练的劳苦。\par

11. 祂所造的一切，在其时令中都是美好的；祂又将整个世界安置在他们心中，以致无人能从始至终查明天主的作为。\par

这是真实的。因为无人能完全理解天主的作为。此外，世界是天主的作为。所以，无人能查明这世界从始至终的时空，即为其指定的时期和预先确定的界限，因为天主将整个世界作为\Emph{一个范围}的无知安置在我们心中。因此有人说：\BibleQuote{求祢向我宣告我日头的短促。}如此，且为我们的益处，这世界（世代）的终结——即今世的生命——是我们所不知道的。\par

\SourceRecord{Translated by S.D.F. Salmond.；Ante-Nicene Fathers；Vol. 6.；Edited by Alexander Roberts, James Donaldson, and A. Cleveland Coxe.；Buffalo, NY: Christian Literature Publishing Co.,；1886.；<http://www.newadvent.org/fathers/0613.htm>.}
